% Chapters/7-Weak-Interactions.tex

\chapter{Weak Interactions}
\label{ch:weak_interactions}

\section{Why do we need weak interactions? (a first phenomenological picture)}
A striking experimental fact is that not all unstable hadrons decay with comparable widths.
Typical widths of hadron resonances (like $\rho,\omega,\Delta$) are of order
\begin{equation}
\Gamma_{\rm strong}\sim 1\text{--}100~\text{MeV}.
\end{equation}
However, some of the \emph{lightest} hadrons in the lowest isospin multiplets (such as $\pi^\pm,\pi^0,p,n$)
have decay widths that are \emph{many orders of magnitude smaller}. The qualitative message is:
\begin{itemize}
  \item Isospin is (approximately) conserved by the \textbf{strong} interactions.
  \item Therefore, the lightest hadrons in isospin multiplets must decay through \textbf{other} interactions.
\end{itemize}

Two classic examples:
\begin{itemize}
  \item $\pi^0$ decays electromagnetically with a width of order $\Gamma_{\rm em}\sim 10^{-5}~\text{MeV}$.
  \item $\pi^\pm$ and the neutron decay weakly, with much smaller widths $\Gamma_{\rm weak}\sim 10^{-14}$--$10^{-24}~\text{MeV}$.
\end{itemize}

Weak interactions are also required to explain quantitatively:
\begin{itemize}
  \item Muon decay:
  \begin{equation}
    \Gamma_{\mu^- \to e^- \bar\nu_e \nu_\mu}\;=\;2.6\times 10^{-16}~\text{MeV}.
  \end{equation}
  \item Decays of the lightest strange hadrons (example numerical widths):
  \begin{align}
    \Gamma_{K^\pm} &= 4.6\times 10^{-16}~\text{MeV}, &
    \Gamma_{K_S^0} &= 6.3\times 10^{-12}~\text{MeV},\\
    \Gamma_{K_L^0} &= 1.1\times 10^{-14}~\text{MeV}, &
    \Gamma_{\Lambda^0} &= 2.1\times 10^{-12}~\text{MeV}.
  \end{align}
\end{itemize}

\section{Historical origin: beta decay and the neutrino}
Weak interactions were historically discovered in nuclear beta decay:
\begin{align}
  {}^{A}_{Z}X &\to {}^{A}_{Z+1}X + e^- \qquad (\beta^-~\text{decay}),\\
  {}^{A}_{Z}X &\to {}^{A}_{Z-1}X + e^+ \qquad (\beta^+~\text{decay}).
\end{align}
In these decays, the observed electron (or positron) energy spectrum is continuous.
To explain the missing energy, momentum, and angular momentum, Pauli postulated the existence
of the neutrino $\nu$: a spin-$\tfrac12$ particle, originally taken to be massless.

At the nucleon level, beta decay can be understood in terms of
\begin{equation}
n \to p\,e^-\,\bar\nu_e,
\qquad
p \to n\,e^+\,\nu_e.
\end{equation}

\subsection{Fermi theory (parity invariant version)}
The first interaction proposed by Fermi was a local four-fermion interaction of the form
\begin{equation}
\mathcal{L}_{\rm int}
=
G\,\bar\psi_p\gamma^\mu\psi_n\;\bar\psi_e\gamma_\mu\psi_{\nu_e}
+\text{H.c.}
\equiv
G\,\bar p\gamma^\mu n\;\bar e\gamma_\mu\nu_e+\text{H.c.}
\end{equation}
This interaction is \textbf{parity invariant}.

However, a parity invariant weak theory runs into difficulties when trying to accommodate
the observed decay pattern of the kaon. Indeed, for $K^+$ decays:
\begin{itemize}
  \item From $K^+ \to \pi^+\pi^0$ one would infer $P(K^+)=+$.
  \item From $K^+ \to \pi^+\pi^0\pi^0$ one would infer $P(K^+)=-$.
\end{itemize}
This tension points to the need for \textbf{parity violation} in weak interactions.

\section{Parity violation and the V--A structure}
\subsection{Lee--Yang proposal and Wu experiment}
Lee and Yang proposed (1956) a direct test of parity violation in beta decay using polarized nuclei.
The idea: polarize ${}^{60}$Co and study the angular distribution of emitted electrons in
\begin{equation}
{}^{60}\text{Co}\to {}^{60}\text{Ni}^\ast\,e^-\,\bar\nu_e.
\end{equation}
The magnetic field used to polarize the nucleus is invariant under parity.
Therefore:
\begin{itemize}
  \item If parity were conserved, the number of electrons emitted along and opposite
        to the polarization direction would be the same.
  \item Wu's experiment (1956) observed a clear asymmetry: electrons preferentially go
        opposite to the polarization direction.
\end{itemize}
This suggests a chiral structure: \textbf{only the left-handed component of leptons
is sensitive to weak interactions.}

\subsection{V--A (vector minus axial) Fermi interaction}
Fermi's original proposal was modified to the \textbf{V--A} form
\begin{equation}
\boxed{
\mathcal{L}_{\rm int}
=
2\sqrt{2}\,G\;\bar p\,\gamma^\mu P_L\,n\;\bar e\,\gamma_\mu P_L\,\nu_e
+\text{H.c.}
\qquad
P_L=\frac{1-\gamma^5}{2}
}
\end{equation}
Key qualitative consequences:
\begin{itemize}
  \item If the outgoing charged lepton mass is negligible, the charged lepton is
        produced \emph{fully polarized}.
  \item All fields are evaluated at the same spacetime point $\Rightarrow$ it is a
        \textbf{zero-range} interaction.
  \item The coupling has mass dimension $[G]=-2$, suggesting an \textbf{effective}
        interaction arising as a low-energy limit of a more fundamental finite-range theory.
\end{itemize}

\subsection{Finite-range origin: massive vector boson mediator}
A zero-range interaction can emerge as the low-energy limit of a massive vector boson exchange.
If $W_\mu$ is a massive vector field with mass $m_W$, then its propagator schematically is
\begin{equation}
\langle 0|T\{W_\mu(x)W_\nu^\dagger(y)\}|0\rangle
=
\int\frac{d^4k}{(2\pi)^4}
\frac{i\,e^{-ik\cdot(x-y)}}{k^2-m_W^2}
\left(-g_{\mu\nu}+\frac{k_\mu k_\nu}{m_W^2}\right).
\end{equation}
In the low-momentum regime $k_\mu\ll m_W$,
\begin{align}
\frac{1}{k^2-m_W^2}
&\simeq
-\frac{1}{m_W^2},\\[1mm]
\Rightarrow\quad
\langle 0|T\{W_\mu(x)W_\nu^\dagger(y)\}|0\rangle
&\simeq
\int\frac{d^4k}{(2\pi)^4}\;e^{-ik\cdot(x-y)}\;\frac{i g_{\mu\nu}}{m_W^2}
=
\frac{i g_{\mu\nu}}{m_W^2}\,\delta^{(4)}(x-y).
\end{align}
Hence one expects parametrically
\begin{equation}
\boxed{G\sim \frac{1}{m_W^2}.}
\end{equation}

\section{Fermi theory at the quark level and global symmetries}
Since $n=(udd)$ and $p=(uud)$, beta decay is naturally interpreted as a quark-level transition
\begin{equation}
d\to u\,e^-\,\bar\nu_e.
\end{equation}

The Fermi theory can be written in terms of quark and lepton fields as
\begin{equation}
\boxed{
\mathcal{L}_{\rm int}
=
2\sqrt{2}\,G\;J^\mu J_\mu^\dagger,
\qquad
J^\mu=\bar u\gamma^\mu P_L d
+\bar\nu_e\gamma^\mu P_L e
+\bar\nu_\mu\gamma^\mu P_L \mu+\cdots
}
\end{equation}
Immediate consequences emphasized by the theory:
\begin{itemize}
  \item If a right-handed neutrino exists, it does not interact in this framework
        (a \textbf{sterile} neutrino).
  \item The interaction is invariant under independent global phase rotations
        of each left-handed lepton doublet:
        \begin{equation}
          \ell_L\to e^{i\theta_\ell}\ell_L,\qquad
          \nu_{\ell L}\to e^{i\theta_\ell}\nu_{\ell L},
          \qquad \ell=e,\mu,\dots
        \end{equation}
  \item With the corresponding transformation of right-handed charged leptons
        $\ell_R\to e^{i\theta_\ell}\ell_R$, the full Lagrangian remains invariant for
        constant $\theta_\ell$ $\Rightarrow$ individual leptonic numbers (electronic, muonic, \dots)
        are conserved.
  \item Similarly, invariance under a global baryon rotation $q\to e^{i\theta_B}q$ (constant $\theta_B$)
        implies \textbf{baryon number conservation}.
\end{itemize}

A key question for low-energy processes is: if the weak interaction is written at the quark level,
how should it be represented at the \textbf{hadronic} level, where nucleons and pions are the degrees
of freedom? The answer involves symmetry constraints (chiral symmetry) and, in general, form factors.
We will see that chiral symmetry strongly constrains the hadronic structure of weak interactions.

\section{Low-energy test I: charged pion decay and helicity suppression}
\subsection{Process and relevant interaction}
Consider the charged pion decay
\begin{equation}
\pi^+ \to \ell^+ \nu_\ell,
\qquad \ell=\mu,e.
\end{equation}
The relevant piece of the Fermi Lagrangian is
\begin{equation}
\mathcal{L}_{\rm int}
=
2\sqrt{2}\,G\;J^{\mu\dagger}_q J_{\ell\mu},
\qquad
J^{\mu\dagger}_q=\bar d\gamma^\mu P_L u,
\qquad
J^\mu_\ell=\bar\nu_\ell\gamma^\mu P_L \ell.
\end{equation}

At first order,
\begin{equation}
\mathcal{M}
=
\langle f|\mathcal{L}_{\rm int}(0)|i\rangle
=
2\sqrt{2}\,G\;
\underbrace{\langle 0|\bar d\gamma^\mu P_L u|\pi^+(p_A)\rangle}_{\text{hadronic}}
\;
\underbrace{\langle \ell^+(p_1,\lambda_1)\,\nu_\ell(p_2,\lambda_2)|\bar\nu_\ell\gamma_\mu P_L \ell|0\rangle}_{\text{leptonic}}.
\end{equation}

Important helicity observation:
\begin{itemize}
  \item The outgoing neutrino is left-handed $\Rightarrow$ its helicity is fixed: $\lambda_2=-1$.
  \item If the charged lepton mass were negligible, the outgoing $\ell^+$ would be produced
        with fixed helicity $\lambda_1=+1$ (fully polarized).
  \item For a pion at rest ($J^P=0^-$), angular momentum conservation would then fail:
        the amplitude must be proportional to $m_\ell$.
\end{itemize}
This is the origin of \textbf{helicity suppression}.

\subsection{Leptonic matrix element}
For the leptonic part,
\begin{equation}
\langle \ell^+(p_1,\lambda_1)\,\nu_\ell(p_2,\lambda_2)|\bar\nu_\ell\gamma^\mu P_L \ell|0\rangle
=
\bar u(2)\gamma^\mu P_L v(1).
\end{equation}

\subsection{Hadronic matrix element and the pion decay constant}
For the hadronic part, define
\begin{equation}
\langle 0|\bar d\gamma^\mu P_L u|\pi^+(p_A)\rangle
\equiv
\frac{i}{\sqrt{2}}\,F_\pi\,p_A^\mu.
\end{equation}
Since the pion has $J^P=0^-$, parity implies that only the axial part contributes:
\begin{equation}
\bar d\gamma^\mu P_L u
=
\frac12\,\bar d\gamma^\mu u
-\frac12\,\bar d\gamma^\mu\gamma^5 u,
\qquad
\Rightarrow\quad
\langle 0|\bar d\gamma^\mu u|\pi^+\rangle=0,
\end{equation}
and
\begin{equation}
\langle 0|\bar d\gamma^\mu P_L u|\pi^+\rangle
=
-\frac12\,\langle 0|\bar d\gamma^\mu\gamma^5 u|\pi^+\rangle.
\end{equation}
The constant $F_\pi$ is the (weak) \textbf{pion decay constant}.

\subsection{Amplitude and explicit $m_\ell$ factor}
Putting both parts together,
\begin{equation}
\mathcal{M}
=
2\sqrt{2}\,G\left(\frac{i}{\sqrt{2}}F_\pi p_A^\mu\right)\bar u(2)\gamma_\mu P_L v(1)
=
2iG F_\pi\,p_A^\mu\,\bar u(2)\gamma_\mu P_L v(1).
\end{equation}
Now use $p_A=p_1+p_2$ and the spinor equations of motion (as used in the slides):
\begin{align}
\bar u(2)\slashed{p}_2 &= 0,\\
(\slashed{p}_1 + m_\ell)\,v(1) &= 0,
\end{align}
so that
\begin{equation}
\bar u(2)\slashed{p}_A P_L v(1)
=
\bar u(2)(\slashed{p}_1+\slashed{p}_2)P_L v(1)
=
\bar u(2)\slashed{p}_1 P_L v(1)
=
- m_\ell\,\bar u(2)P_R v(1).
\end{equation}
Therefore,
\begin{equation}
\boxed{
\mathcal{M}=-2iG F_\pi\,m_\ell\;\bar u(2)P_R v(1)
}
\end{equation}
This makes helicity suppression manifest: $\mathcal{M}\propto m_\ell$.

\subsection{Polarization selection and decay width}
In the pion rest frame, the slides compute explicitly
\begin{align}
\mathcal{M}\!\left(\pi^+(p_A=0)\to \ell^+(p_1,+)\,\nu_\ell(p_2,-)\right) &= 0,\\
\mathcal{M}\!\left(\pi^+(p_A=0)\to \ell^+(p_1,-)\,\nu_\ell(p_2,-)\right) &\neq 0,
\end{align}
i.e.\ the charged lepton is fully polarized in this approximation.

Summing over helicities (as done in the slides) yields
\begin{equation}
\sum_{\lambda_1=\pm}|\mathcal{M}|^2
=
8G^2F_\pi^2 m_\ell^2\,E_2\,m_\pi,
\end{equation}
with $p=E_2=\sqrt{E_1^2-m_\ell^2}$ and $E_1+E_2=m_\pi$ in the pion rest frame.
Using
\begin{equation}
p=\frac{m_\pi^2-m_\ell^2}{2m_\pi},
\end{equation}
the decay width becomes
\begin{equation}
\boxed{
\Gamma(\pi^+\to \ell^+\nu_\ell)
=
\frac{G^2F_\pi^2 m_\ell^2 m_\pi}{4\pi}
\left(1-\frac{m_\ell^2}{m_\pi^2}\right)^2
}
\end{equation}

\subsection{A parameter-free prediction: the $e/\mu$ ratio}
Taking the ratio cancels $G$ and $F_\pi$:
\begin{equation}
\boxed{
\frac{\Gamma(\pi^+\to e^+\nu_e)}{\Gamma(\pi^+\to \mu^+\nu_\mu)}
=
\frac{m_e^2}{m_\mu^2}
\left(\frac{m_\pi^2-m_e^2}{m_\pi^2-m_\mu^2}\right)^2
\;\simeq\;1.2\times 10^{-4},
}
\end{equation}
in agreement with the quoted experimental value $\simeq 1.2\times 10^{-4}$.
This is a clean low-energy confirmation of the V--A structure.

\section{Non-linear sigma model: implementing weak interactions at hadronic level}
Pion decay and beta decay are low-energy processes, so they should be describable
within the non-linear sigma model (effective hadronic theory).
The guiding principle is to \textbf{implement the symmetry breaking pattern}.

\subsection{External left-handed field and local chiral symmetry}
Start from the quark-level charged current:
\begin{equation}
\mathcal{L}_{\rm int}=2\sqrt{2}\,G\;J^{\mu\dagger}_q J_{\ell\mu}+\text{H.c.},
\qquad
J^{\mu\dagger}_q=\bar d\gamma^\mu P_L u.
\end{equation}
Introduce the isospin doublet $q=(u,d)^T$ and
\begin{equation}
\tau^+ =
\begin{pmatrix}
0&1\\
0&0
\end{pmatrix},
\qquad
\tau^- =
\begin{pmatrix}
0&0\\
1&0
\end{pmatrix},
\end{equation}
so
\begin{equation}
J_q^{\mu\dagger}=\bar q\gamma^\mu P_L \tau^- q,
\qquad
J_q^\mu=\bar q\gamma^\mu P_L \tau^+ q.
\end{equation}
Then one can write
\begin{equation}
\mathcal{L}_{\rm int}=\bar q\,\gamma^\mu P_L\,a^L_\mu\,q,
\end{equation}
with an \emph{external} matrix field
\begin{equation}
a^L_\mu \equiv 2\sqrt{2}\,G
\begin{pmatrix}
0 & \bar\ell\gamma_\mu P_L \nu_\ell\\
\bar\nu_\ell\gamma_\mu P_L \ell & 0
\end{pmatrix}.
\end{equation}

In the massless limit $m_u=m_d=0$, the QCD Lagrangian has a global
$SU_L(2)\otimes SU_R(2)$ chiral symmetry. The term above breaks it because
$[g_L,a^L_\mu]\neq 0$ for constant $g_L\in SU(2)$. The trick is:
\begin{itemize}
  \item Promote $g_L$ to a \textbf{local} transformation $g_L(x)$ and assign
        the transformation law
        \begin{equation}
        a^L_\mu(x)\to g_L(x)a^L_\mu(x)g_L^\dagger(x)+i g_L(x)\partial_\mu g_L^\dagger(x),
        \qquad g_R=\text{constant}.
        \end{equation}
  \item Then the left-handed quark kinetic term becomes locally invariant if we replace
        $\partial_\mu$ by $(\partial_\mu-i a^L_\mu)$ on $q_L$:
        \begin{equation}
          \mathcal{L}=\bar q_L(i\slashed{\partial}+\slashed{a}_L)q_L+\bar q_R i\slashed{\partial} q_R,
        \end{equation}
        and for QCD the same holds with $\partial\to D$ (including gluons).
\end{itemize}

\subsection{Coupling to the non-linear sigma model}
In the non-linear sigma model, the pion field is encoded in
\begin{equation}
U(x)\in SU(2),
\qquad
U(x)\to g_L(x)U(x)g_R^\dagger.
\end{equation}
Because $g_L$ is local, $\partial_\mu U$ does not transform covariantly.
Define the covariant derivative
\begin{equation}
D_\mu U \equiv (\partial_\mu-i a^L_\mu)\,U,
\qquad
D_\mu U \to g_L(x) (D_\mu U) g_R^\dagger,
\end{equation}
and replace
\begin{equation}
\mathcal{L}=\frac{f_\pi^2}{4}\,\mathrm{tr}(\partial_\mu U^\dagger \partial^\mu U)
\quad\longrightarrow\quad
\mathcal{L}=\frac{f_\pi^2}{4}\,\mathrm{tr}\!\left[(D_\mu U)^\dagger D^\mu U\right].
\end{equation}

Expanding in $G$ gives the weak interaction term (to $\mathcal{O}(G)$):
\begin{align}
\mathcal{L}^{\rm weak}_{\rm int}
&=
\frac{f_\pi^2}{4}\,\mathrm{tr}\!\Big(\partial_\mu U^\dagger(-i a_L^\mu)U
+U^\dagger i a_L^\mu \partial_\mu U\Big)+\mathcal{O}(G^2)\\[1mm]
&\simeq
i\,2\sqrt{2}\,G\,\frac{f_\pi^2}{4}\,
\mathrm{tr}\!\left[
\begin{pmatrix}
0 & \bar\ell\gamma_\mu P_L \nu_\ell\\
\bar\nu_\ell\gamma_\mu P_L \ell & 0
\end{pmatrix}
\left(\partial^\mu U\,U^\dagger-U\partial^\mu U^\dagger\right)
\right].
\end{align}
To leading order in pion fields,
\begin{equation}
\boxed{
\mathcal{L}^{\rm weak}_{\rm int}
\simeq
-2G f_\pi\left(
\partial_\mu\pi^+\;\bar\nu_\ell\gamma^\mu P_L \ell
+\partial_\mu\pi^-\;\bar\ell\gamma^\mu P_L \nu_\ell
\right).
}
\end{equation}
Comparing with the quark-level interaction identifies the hadronic representation
of the current:
\begin{equation}
J_q^{\mu\dagger}
\;\to\;
\frac{i f_\pi^2}{4}\,
\mathrm{tr}\!\left[
\tau^-\left(\partial^\mu U\,U^\dagger-U\partial^\mu U^\dagger\right)
\right]
\simeq
-\frac{f_\pi}{\sqrt{2}}\,\partial^\mu \pi^+ + \cdots
\end{equation}

\subsection{Recovering pion decay and the relation $F_\pi=f_\pi$}
Using the hadronic current above,
\begin{equation}
\langle 0|\bar d\gamma^\mu P_L u|\pi^+(p_A)\rangle
\;\parallel\;
\left\langle 0\left|-\frac{f_\pi}{\sqrt2}\,\partial^\mu\pi^+\right|\pi^+(p_A)\right\rangle
=
\frac{i f_\pi}{\sqrt2}\,p_A^\mu,
\end{equation}
so we conclude
\begin{equation}
\boxed{F_\pi=f_\pi.}
\end{equation}
This is a powerful EFT statement: $f_\pi$ not only controls low-energy \textbf{strong} pion interactions,
it also fixes the normalization of low-energy \textbf{weak} pion processes.
Moreover, at low energy one can compute other weak pion processes (e.g.\ $\pi^+\to \pi^0 e^+\nu_e$)
by expanding the same current to higher order in pion fields, without introducing new parameters.

\section{Weak interactions with nucleons and the axial coupling $g_A$}
For nucleons, the non-linear sigma model (as used in the slides) contains
\begin{equation}
\bar N\left(
i\slashed{\partial}+\slashed{v}+i g_A \slashed{a}\gamma^5 - m_N
\right)N,
\qquad
v_\mu=\frac12\left(u\partial_\mu u^\dagger+u^\dagger\partial_\mu u\right),
\quad
a_\mu=\frac12\left(u\partial_\mu u^\dagger-u^\dagger\partial_\mu u\right),
\quad U=u^2.
\end{equation}
To incorporate the weak external field $a^L_\mu(x)$ and local $SU_L(2)$, one replaces
\begin{equation}
\partial_\mu u \to D_\mu u=(\partial_\mu-i a^L_\mu)u,
\end{equation}
and correspondingly
\begin{equation}
v_\mu\to v'_\mu\equiv \frac12\left(u\partial_\mu u^\dagger+u^\dagger D_\mu u\right),
\qquad
a_\mu\to a'_\mu\equiv \frac12\left(u\partial_\mu u^\dagger-u^\dagger D_\mu u\right).
\end{equation}
This yields, at $\mathcal{O}(G)$, the nucleon weak interaction term
\begin{equation}
\mathcal{L}^{\rm weak}_{\rm int}
=
\frac12\,\bar N\left(u^\dagger \slashed{a}_L u - g_A\,u^\dagger \slashed{a}_L u\,\gamma^5\right)N.
\end{equation}

\subsection{Hadronic charged current at nucleon level}
Comparing with the quark-level interaction gives the hadronic representation
\begin{equation}
J_q^{\mu\dagger}
\;\to\;
\frac12\,\bar N\gamma^\mu\left(u^\dagger\tau^- u - g_A\,u^\dagger\tau^- u\,\gamma^5\right)N
\simeq
\bar n\,\gamma^\mu\frac{1-g_A\gamma^5}{2}\,p+\cdots
\end{equation}
For processes with \emph{no pions} (like ordinary beta decay) we set $u=1$ and obtain
\begin{equation}
\boxed{
\mathcal{L}^{\rm weak}_{\rm int}
=
2\sqrt{2}\,G\;
\bar n\,\gamma^\mu\frac{1-g_A\gamma^5}{2}\,p\;
\bar\nu_\ell\gamma_\mu P_L \ell
+\text{H.c.}
}
\end{equation}
A crucial point emphasized in the slides:
\begin{equation}
g_A\simeq 1.27
\quad\Rightarrow\quad
\text{at nucleon level the weak interaction is \emph{not purely left-handed}.}
\end{equation}
Indeed,
\begin{equation}
\frac{1-g_A\gamma^5}{2}
=
\frac{1+g_A}{2}\,P_L+\frac{1-g_A}{2}\,P_R
\approx
1.13\,P_L+0.13\,P_R.
\end{equation}

\section{Nuclear beta decay as a low-energy limit}
Nuclei are bound states of non-relativistic nucleons. Hence nuclear beta decay is a
low-energy process, and we can start from the nucleon-level interaction and further
simplify it using non-relativistic (Schr\"odinger) fields.

\subsection{Non-relativistic reduction and Fermi vs.\ Gamow--Teller}
For $|\vec p|\ll m$,
\begin{equation}
u_\lambda(\vec p)=\sqrt{E+m}
\begin{pmatrix}
\chi_\lambda\\
\frac{\vec\sigma\cdot\vec p}{E+m}\chi_\lambda
\end{pmatrix}
\simeq
\sqrt{2m}
\begin{pmatrix}
\chi_\lambda\\
0
\end{pmatrix}.
\end{equation}
Write the nucleon field as
\begin{equation}
N(x)=e^{-imx^0}P_+\,\phi_N(x),
\qquad
P_+=\frac{1+\gamma^0}{2}=
\begin{pmatrix}
1&0\\
0&0
\end{pmatrix},
\end{equation}
where $\phi_N(x)$ is a Schr\"odinger spin-$\tfrac12$ field. Using the projector identities quoted:
\begin{equation}
P_+\gamma^\mu P_+=(1,\vec 0),
\qquad
P_+\gamma^\mu\gamma^5 P_+=(0,\vec\sigma),
\end{equation}
the interaction becomes
\begin{equation}
\boxed{
\mathcal{L}^{\rm weak}_{\rm int}
=
\sqrt{2}\,G\left(
\phi_n^\dagger \phi_p\;\bar\nu_\ell\gamma^0 P_L \ell
- g_A\,\phi_n^\dagger\sigma^i\phi_p\;\bar\nu_\ell\gamma^i P_L \ell
\right)e^{-i(m_p-m_n)x^0}
+\text{H.c.}
}
\end{equation}
Interpretation:
\begin{itemize}
  \item The spin-independent term $\phi_n^\dagger\phi_p$ induces \textbf{Fermi transitions}.
  \item The spin-dependent term $\phi_n^\dagger\sigma^i\phi_p$ induces \textbf{Gamow--Teller transitions}.
  \item If initial and final nuclear states have $J^P=0^+$, the Gamow--Teller term does not contribute,
        so the decay width does not depend on $g_A$.
\end{itemize}

\subsection{Isospin matrix elements and accurate determination of $G$}
If, in addition, the initial and final nuclear states belong to the same isospin multiplet,
then the form factor
\begin{equation}
\langle {}^{A}_{Z-1}X|\phi_n^\dagger(0)\phi_p(0)|{}^{A}_{Z}X\rangle
\end{equation}
can be evaluated exactly, because $\phi_n^\dagger\phi_p$ replaces a proton by a neutron
and plays the role of a $\tau^-$ operator.

Example: ${}^{14}_6\text{C}$, ${}^{14}_7\text{N}^\ast$, ${}^{14}_8\text{O}$ have $J^P=0^+$ and form an isospin triplet.
For ${}^{14}_8\text{O}\to {}^{14}_7\text{N}^\ast\,e^+\,\nu_e$,
\begin{equation}
\langle {}^{14}_7\text{N}^\ast|\phi_n^\dagger(0)\phi_p(0)|{}^{14}_8\text{O}\rangle
\to
\langle I I_3'|\tau^-|I I_3\rangle
=
\sqrt{2}.
\end{equation}
In such cases, $G$ can be measured accurately:
\begin{equation}
\boxed{G=1.136(3)\times 10^{-5}~\text{GeV}^{-2}\qquad(\beta\text{-decay determination}).}
\end{equation}

\section{Example calculation: spectrum and total width in nuclear beta decay}
Consider (as in the slides)
\begin{equation}
{}^{14}_8\text{O}_A \to {}^{14}_7\text{N}^\ast_{1} + e^+_2 + \nu_e{}_3.
\end{equation}
Using the Fermi term only, the amplitude is
\begin{equation}
\mathcal{M}
=
\sqrt{2}\,G\;\langle 1|\phi_n^\dagger\phi_p|A\rangle\;
\langle 2,3|\bar\nu_\ell\gamma^0 P_L \ell|0\rangle.
\end{equation}
Inserting the isospin matrix element $\sqrt2$ and the relativistic normalizations used in the slides,
\begin{equation}
\mathcal{M}
=
\sqrt{2}\,G
\left(\sqrt2\right)
\left(\sqrt{2m_A}\sqrt{2m_1}\right)
\bar u(3)\gamma^0 P_L v(2).
\end{equation}

If the positron polarization is not measured, one obtains (trace steps as in the slides)
\begin{align}
|\mathcal{M}|^2
&=
16G^2 m_A m_1
\sum_{\lambda_2,\lambda_3}
\mathrm{tr}\!\left[
\slashed{p}_3\gamma^0 P_L (\slashed{p}_2-m_2)\gamma^0 P_L
\right]\\
&=
32G^2 m_A m_1\left(E_2E_3+|\vec p_2||\vec p_3|\cos\theta\right),
\qquad
\hat p_2\cdot\hat p_3\equiv \cos\theta.
\end{align}

\subsection{Phase space and spectrum (massless neutrino, $m_e$ neglected)}
The decay width is
\begin{equation}
\Gamma
=
\frac{1}{2m_A}
\int\frac{d^3\vec p_1}{(2\pi)^3 2E_1}
\int\frac{d^3\vec p_2}{(2\pi)^3 2E_2}
\int\frac{d^3\vec p_3}{(2\pi)^3 2E_3}\;
|\mathcal{M}|^2\;
(2\pi)^4\delta(m_A-E_1-E_2-E_3)\delta^{(3)}(\vec p_1+\vec p_2+\vec p_3).
\end{equation}
Carrying out the integrals as done in the slides and using $E_1\simeq m_1$, one arrives at
\begin{equation}
\frac{d\Gamma}{dE_2}
=
\frac{G^2}{\pi^3}\,|\vec p_2|E_2\,(E_0-E_2)^2
\simeq
\frac{G^2}{\pi^3}\,E_2^2\,(E_0-E_2)^2
\qquad (E_2\gg m_e),
\end{equation}
where
\begin{equation}
E_0 \equiv m_A-m_1.
\end{equation}
The total width becomes
\begin{equation}
\boxed{
\Gamma
\simeq
\int_0^{E_0}dE_2\;\frac{G^2}{\pi^3}E_2^2(E_0-E_2)^2
=
\frac{G^2 E_0^5}{30\pi^3}.
}
\end{equation}

\subsection{Remarks (as emphasized in the slides)}
\begin{itemize}
  \item The width is extremely sensitive to the nuclear mass difference:
        \begin{equation}
        \Gamma\sim E_0^5=(m_A-m_1)^5.
        \end{equation}
  \item The energy spectrum is symmetric around $E_2\simeq E_0/2$ and peaks there.
  \item The electron/positron mass is not always negligible (especially near the endpoint).
  \item Coulomb interaction between the charged lepton and the nucleus affects fine details of the spectrum
        and the accurate decay width.
\end{itemize}

\section{Finite neutrino mass effects on beta decay endpoint}
Neutrinos are massless in the Standard Model, but oscillations imply a small neutrino mass.
If a finite neutrino mass $m_\nu$ is included:
\begin{itemize}
  \item The amplitude remains exactly the same.
  \item The phase space is modified. The decay width involves
        \begin{equation}
        \sqrt{E_3^2-m_\nu^2},
        \qquad E_3=E_0-E_2,
        \end{equation}
        and the endpoint shifts to $E_2^{\rm max}=E_0-m_\nu$.
\end{itemize}
In the approximation used in the slides ($E_2\gg m_e$, $E_0\gg m_e$),
\begin{equation}
\boxed{
\frac{d\Gamma}{dE_2}
\simeq
\frac{G^2}{\pi^3}\,E_2^2\,(E_0-E_2)\,\sqrt{(E_0-E_2)^2-m_\nu^2}.
}
\end{equation}
Near the endpoint, $E_2=E_0-m_\nu-\delta E_2$ with $\delta E_2\to 0$,
\begin{equation}
\frac{d\Gamma}{dE_2}
\simeq
\frac{G^2}{\pi^3}\,E_0^2\,m_\nu\,\sqrt{2m_\nu\,\delta E_2},
\end{equation}
which highlights how endpoint measurements are sensitive to $m_\nu$.

\section{Muon decay: determination of $G$ and comparison with beta decay}
Consider
\begin{equation}
\mu^-_A \to e^-_1\,\bar\nu_e{}_2\,\nu_\mu{}_3.
\end{equation}
Since $m_\mu\simeq 106~\text{MeV}\gg m_e\simeq 0.5~\text{MeV}$, neglect $m_e$.
The relevant interaction is
\begin{equation}
\mathcal{L}_{\rm int}
=
2\sqrt{2}\,G\;J_e^{\mu\dagger}J_{\mu\mu}+\text{H.c.},
\qquad
J_e^{\mu\dagger}=\bar e\gamma^\mu P_L \nu_e,
\qquad
J_{\mu}^\mu=\bar\nu_\mu\gamma^\mu P_L \mu.
\end{equation}
The amplitude is
\begin{equation}
\mathcal{M}
=
2\sqrt{2}\,G\;\bar u(3)\gamma^\mu P_L u(A)\;\bar u(1)\gamma_\mu P_L v(2).
\end{equation}

If the muon is unpolarized and final polarizations are not measured, the slides show
\begin{align}
|\mathcal{M}|^2
&=
4G^2\;
\mathrm{tr}\!\Big(\slashed{p}_3\gamma^\mu P_L(\slashed{p}_A+m_A)\gamma^\nu P_L\Big)\;
\mathrm{tr}\!\Big(\slashed{p}_1\gamma_\mu P_L\slashed{p}_2\gamma_\nu P_L\Big)\\
&\equiv
4G^2\,L^{\mu\nu}(p_3,p_A)\,L_{\mu\nu}(p_1,p_2)
=
64G^2\,(p_A\!\cdot p_2)(p_3\!\cdot p_1).
\end{align}
Using $(p_A\!\cdot p_2)=m_A E_2$ and $(p_3\!\cdot p_1)=\frac12(p_A-p_2)^2=\frac{m_A^2}{2}-m_AE_2$,
the amplitude depends only on energies (no angles), simplifying the phase-space integration.

The electron energy spectrum (with $m_A=m_\mu$) is
\begin{equation}
\boxed{
\frac{d\Gamma}{dE_1}
=
\frac{G^2 m_\mu^2 E_1^2}{4\pi^3}\left(1-\frac{4E_1}{3m_\mu}\right).
}
\end{equation}
The total muon decay width is
\begin{equation}
\boxed{
\Gamma_\mu
=
\frac{G^2 m_\mu^5}{192\pi^3}.
}
\end{equation}
This allows an independent measurement:
\begin{equation}
G\big|_{\mu\text{-decay}}=1.16632(2)\times 10^{-5}~\text{GeV}^{-2}.
\end{equation}
The slides emphasize a small but significant discrepancy with the beta-decay value
$G\big|_{\beta\text{-decay}}=1.136(3)\times 10^{-5}~\text{GeV}^{-2}$, which motivates the next section.

\section{Strange quarks, kaon decays, and Cabibbo mixing}
\subsection{Failure of universal $d\to u$ coupling for $s\to u$}
Consider $K^+\to e^+\nu_e$. The slides quote
\begin{equation}
\Gamma(K^+\to e^+\nu_e)=1.28\times 10^3~\text{s}^{-1}.
\end{equation}
If the $s$ quark coupled to $u$ with the same strength as $d$ does, then (neglecting $m_e$)
\begin{equation}
\frac{\Gamma(K^+\to e^+\nu_e)}{\Gamma(\pi^+\to e^+\nu_e)}
\simeq
\frac{m_K f_K^2}{m_\pi f_\pi^2}.
\end{equation}
Approximate chiral $SU_L(3)\times SU_R(3)$ suggests $f_K\simeq f_\pi$, leading to the rough expectation
\begin{equation}
\left.\frac{\Gamma(K^+\to e^+\nu_e)}{\Gamma(\pi^+\to e^+\nu_e)}\right|_{\rm th}
\simeq
\frac{m_K}{m_\pi}\simeq 3.5.
\end{equation}
But experimentally the ratio is $\simeq 0.27$. Conclusion:
\begin{equation}
\boxed{\text{The $s$ quark does not couple to $u$ with the same strength as $d$ does.}}
\end{equation}

\subsection{Introducing $V_{ud}$ and $V_{us}$ (Cabibbo idea)}
Modify the current by introducing parameters $V_{ud}$ and $V_{us}$:
\begin{equation}
J^\mu
=
\bar\nu_e\gamma^\mu P_L e
+\bar\nu_\mu\gamma^\mu P_L \mu
+V_{ud}\,\bar u\gamma^\mu P_L d
+V_{us}\,\bar u\gamma^\mu P_L s+\cdots
\end{equation}
Assuming $V_{ud}\simeq 1$, the kaon/pion ratio suggests $|V_{us}|^2\sim 0.27$.
A more accurate determination comes from $K^0\to \pi^- e^+\bar\nu_e$ (not depending on $f_\pi$ or $f_K$):
\begin{equation}
|V_{us}|\simeq 0.2249(10).
\end{equation}

Cabibbo parameterized
\begin{equation}
\sin\theta_c=V_{us},
\qquad
\cos\theta_c=V_{ud},
\end{equation}
motivated by assuming that the quark fields appearing in the weak current are unitary rotations
of the mass eigenstate fields:
\begin{equation}
\begin{pmatrix}
d'\\
s'
\end{pmatrix}
=
\begin{pmatrix}
\cos\theta_c & \sin\theta_c\\
-\sin\theta_c & \cos\theta_c
\end{pmatrix}
\begin{pmatrix}
d\\
s
\end{pmatrix},
\qquad
J^\mu=\bar\nu_e\gamma^\mu P_L e+\bar\nu_\mu\gamma^\mu P_L\mu+\bar u\gamma^\mu P_L d'+\cdots
\end{equation}

\subsection{More quarks and resolving the $G$ discrepancy}
A natural question raised in the slides: if only $d'$ couples to $u$, who couples to $s'$?
This motivates introducing a new quark $c$ (charm) that couples to $s'$:
\begin{equation}
J^\mu=\bar\nu_e\gamma^\mu P_L e+\bar\nu_\mu\gamma^\mu P_L\mu+\bar u\gamma^\mu P_L d'+\bar c\gamma^\mu P_L s'+\cdots
\end{equation}

With Cabibbo's modification, the beta-decay determination is consistent with muon decay:
\begin{equation}
G_{\beta}\;=\;1.136(3)\times 10^{-5}~\text{GeV}^{-2}
\;\;\parallel\;\;
G_{\mu}\cos\theta_c
=
1.16632(2)\times 10^{-5}~\text{GeV}^{-2}\times 0.9744
=
1.13644\times 10^{-5}~\text{GeV}^{-2}.
\end{equation}
So the discrepancy is resolved by the Cabibbo factor.

\section{Counting parameters in quark mixing matrices: towards CKM}
For $N$ down-type quarks, the most general mixing is an $N\times N$ unitary matrix $V$ with
\begin{equation}
\sum_i V_{ij}V^\ast_{ik}=\delta_{jk}.
\end{equation}
The slides count the independent parameters:
\begin{itemize}
  \item An arbitrary complex $N\times N$ matrix has $N^2$ moduli and $N^2$ phases.
  \item Unitarity imposes constraints:
        \begin{itemize}
          \item $N$ constraints for $j=k$ on moduli,
          \item $N(N-1)/2$ constraints for $j\neq k$ on both moduli and phases.
        \end{itemize}
  \item After imposing unitarity, the remaining are
        \begin{equation}
        \text{Moduli: } \frac{N(N-1)}{2},
        \qquad
        \text{Phases: } \frac{N(N+1)}{2}.
        \end{equation}
  \item Additionally, $2N-1$ phases can be absorbed by redefining quark fields.
\end{itemize}
Final count:
\begin{equation}
\boxed{
\text{Independent moduli: }\frac{N(N-1)}{2},
\qquad
\text{Independent phases: }\frac{(N-2)(N-1)}{2}.
}
\end{equation}
Checks:
\begin{itemize}
  \item $N=2$: one modulus, no phase (Cabibbo angle).
  \item $N=3$: three moduli and \textbf{one phase} $\Rightarrow$ CP violation (Kobayashi--Maskawa).
\end{itemize}
For $N=3$ the matrix is the \textbf{CKM matrix}.

\section{CP violation from a complex CKM phase}
For $N=3$ families, a phase in the CKM matrix implies complex terms in the Fermi Lagrangian.
Example term (if $V_{ub}$ is complex):
\begin{equation}
V_{ub}V^\ast_{cs}\;\bar u\gamma^\mu P_L b\;\bar s\gamma_\mu P_L c
+
V^\ast_{ub}V_{cs}\;\bar b\gamma^\mu P_L u\;\bar c\gamma_\mu P_L s.
\end{equation}
Under CP, one finds that the left current transforms into the right current (and vice versa),
but the Lagrangian is invariant only if
\begin{equation}
V_{ub}V^\ast_{cs}=V^\ast_{ub}V_{cs},
\end{equation}
which fails if $V_{ub}$ is complex. Hence CP is violated.

The slides highlight the historical narrative:
\begin{itemize}
  \item CP violation was observed in neutral kaon decays since the 1960s.
  \item The idea that a third family could explain CP violation was realized before the $b$ quark discovery.
  \item The discovery of $b$ (1977) opened the possibility of explaining kaon CP violation in the electroweak framework.
\end{itemize}

\section{Neutral kaons and CP violation (classic system)}
Neutral kaons:
\begin{equation}
K^0=(\bar s d),\qquad \bar K^0=(d\bar s),
\qquad J^P=0^-.
\end{equation}
Define CP eigenstates (using $C|\bar K^0\rangle=|K^0\rangle$ and $C|K^0\rangle=|\bar K^0\rangle$):
\begin{align}
|K_S\rangle &\equiv \frac{1}{\sqrt2}\left(|K^0\rangle-|\bar K^0\rangle\right),
&
|K_L\rangle &\equiv \frac{1}{\sqrt2}\left(|K^0\rangle+|\bar K^0\rangle\right),
\end{align}
with
\begin{equation}
CP|K_S\rangle=|K_S\rangle,\qquad CP|K_L\rangle=-|K_L\rangle.
\end{equation}

Possible final states include $\pi\pi$ and $\pi\pi\pi$:
\begin{itemize}
  \item Two-pion decay: the orbital angular momentum must be $L=0$, and $CP|\pi\pi\rangle=+|\pi\pi\rangle$.
  \item Three-pion decay dominated by $L=0$ gives $CP|\pi\pi\pi\rangle=-|\pi\pi\pi\rangle$.
\end{itemize}
Therefore, if CP were conserved:
\begin{equation}
|K_S\rangle \to |\pi\pi\rangle,
\qquad
|K_L\rangle \to |\pi\pi\pi\rangle.
\end{equation}

Lifetimes:
\begin{equation}
\tau(K_S)\sim 8.9\times 10^{-11}\text{ s},
\qquad
\tau(K_L)\sim 5.1\times 10^{-8}\text{ s}.
\end{equation}
In a long-lived beam, essentially only $K_L$ survive. Experimentally it was observed that
$K_L$ also decays to $\pi\pi$ with branching fractions
\begin{equation}
\mathrm{BF}(K_L\to \pi^+\pi^-)\sim 2\times 10^{-3},
\qquad
\mathrm{BF}(K_L\to \pi^0\pi^0)\sim 8\times 10^{-4},
\end{equation}
which implies
\begin{equation}
\boxed{\text{CP is violated.}}
\end{equation}
Consequences include particle--antiparticle asymmetries such as
\begin{equation}
\Gamma(K_L\to \pi^+ e^- \bar\nu_e)\neq \Gamma(K_L\to \pi^- e^+ \nu_e).
\end{equation}
The slides also stress the cosmological relevance: CP violation is one of the three
Sakharov conditions for explaining the matter--antimatter asymmetry.

\section{T violation and CPT}
The slides state:
\begin{itemize}
  \item The same terms that violate CP in the Fermi Lagrangian also violate T (time reversal).
  \item Direct T violation has been observed (e.g.\ 1998 in $K^0$ decays, CPLEAR; 2012 in $B$ decays, BaBar).
  \item CPT is always respected.
\end{itemize}

\section{CKM matrix hierarchy and size of CP violation}
The CKM matrix is known rather accurately. Its parametric structure is summarized as
\begin{equation}
V_{ij}=
\begin{pmatrix}
\mathcal{O}(1) & \mathcal{O}(\lambda) & \mathcal{O}(\lambda^3)\\
\mathcal{O}(\lambda) & \mathcal{O}(1) & \mathcal{O}(\lambda^2)\\
\mathcal{O}(\lambda^3) & \mathcal{O}(\lambda^2) & \mathcal{O}(1)
\end{pmatrix},
\qquad
\lambda\simeq 0.22,
\end{equation}
with $j=(d,s,b)$ and $i=(u,c,t)$. Interpretation:
\begin{itemize}
  \item If kinematically allowed, quarks prefer to decay within the same family.
  \item Next preference is to the closest family (controlled by powers of $\lambda$).
\end{itemize}
The size of CP violation can be estimated by Jarlskog's determinant,
\begin{equation}
\boxed{\text{Jarlskog invariant}~\sim 2.96\times 10^{-5}},
\end{equation}
while $|\det V|\sim 1$.

\section{Neutrino masses, mixing, and oscillations (PMNS matrix)}
The slides emphasize that the Standard Model has massless neutrinos, but oscillations
observed in nature imply small neutrino masses.

Introducing a right-handed neutrino (which is sterile in the weak current) allows a Dirac mass term.
Then, similarly to quarks, the neutrino fields appearing in the weak currents may not coincide with mass eigenstates,
but be related by a unitary transformation. This leads to the leptonic analogue of CKM: the
\textbf{Pontecorvo--Maki--Nakagawa--Sakata} (PMNS) matrix and to neutrino oscillations.

Observed oscillation evidence (as listed in the slides):
\begin{itemize}
  \item Solar neutrinos (deficit observed since late 1960s; confirmed definitively by SNO in 2001).
  \item Atmospheric neutrinos (imbalance confirmed by Kamiokande in 1998).
  \item Reactor neutrinos.
  \item Beam neutrinos (accelerator-produced).
\end{itemize}

The PMNS matrix is summarized in the slides approximately as
\begin{equation}
U_{\rm PMNS}\approx
\begin{pmatrix}
\sqrt{\frac{2}{3}} & \sqrt{\frac{1}{3}} & \mathcal{O}(\lambda)\\
-\sqrt{\frac{1}{6}} & \sqrt{\frac{1}{3}} & \sqrt{\frac{1}{2}}\\
\sqrt{\frac{1}{6}} & -\sqrt{\frac{1}{3}} & \sqrt{\frac{1}{2}}
\end{pmatrix},
\qquad
\lambda\simeq 0.22,
\end{equation}
with a pattern very different from the quark sector, providing an additional possible
source of CP violation.

If neutrino masses are Majorana:
\begin{itemize}
  \item Two more phases appear $\Rightarrow$ more sources of CP violation.
  \item One does not need a right-handed neutrino.
  \item Lepton number violation should appear, e.g.\ neutrinoless double beta decay
        (example in the slides: ${}^{128}_{52}\text{Te}\to {}^{128}_{54}\text{Xe}\,e^-e^-$).
\end{itemize}

\section{Neutral currents: necessity, suppression of FCNC, and discovery}
\subsection{From charged currents to neutral currents}
So far we have used a charged current $J^\mu$ that changes electric charge by one unit ($\Delta Q=+1$),
and its conjugate $J^{\mu\dagger}$ with $\Delta Q=-1$:
\begin{equation}
J^\mu_{\rm cc}=\sum_F \bar F\,\tau^+\gamma^\mu P_L F',
\qquad
J^{\mu\dagger}_{\rm cc}=\sum_F \bar F'\,\tau^-\gamma^\mu P_L F,
\end{equation}
where $F$ runs over $SU(2)$ doublets such as
\begin{equation}
\begin{pmatrix}\nu_e\\ e\end{pmatrix},
\begin{pmatrix}u\\ d\end{pmatrix},
\begin{pmatrix}\nu_\mu\\ \mu\end{pmatrix},
\begin{pmatrix}c\\ s\end{pmatrix},
\dots
\end{equation}
and $F'=F$ for leptons, while for quarks $F'=V F$ with $V$ the CKM matrix.

If an underlying $SU(2)$ structure exists, then besides $\tau^\pm$ one must also have $\tau^3$,
which leads to \textbf{neutral currents}.

\subsection{Flavor-changing neutral currents (FCNC) are suppressed}
Experimentally, neutral currents that change flavor are very suppressed, e.g.
\begin{equation}
\frac{\Gamma(K_L^0\to \mu^+\mu^-)}{\Gamma(K_L^0\to \text{all})}\simeq 9\times 10^{-9}.
\end{equation}
Therefore we only include neutral currents that do \emph{not} change flavor.
This is natural if CKM is unitary:
\begin{equation}
J^\mu_{\rm nc}\sim \sum_{F'}\bar F'\tau^3\gamma^\mu P_L F'
=
\sum_F \bar F\tau^3\gamma^\mu P_L F.
\end{equation}
(While the full electroweak structure is more complicated, the key feature is invariance under
unitary field redefinitions.)

\subsection{Experimental discovery: Gargamelle (CERN, 1973)}
Neutral currents were suggested (1958, Bludman) and incorporated in early electroweak ideas (Glashow, 1961).
They do not change flavor, hence they compete with electromagnetic interactions and are difficult to observe.

They were first detected at CERN in 1973 by the Gargamelle bubble chamber experiment.
A (anti)neutrino beam was produced from decays $\pi^\pm,K^\pm\to \mu^\pm\nu$,
with parent mesons produced by proton collisions on a Be target. Observed processes:
\begin{align}
\bar\nu_\mu e^- &\to \bar\nu_\mu e^-,\\
\bar\nu_\mu N &\to \bar\nu_\mu X,\qquad (X~\text{hadrons, no muon}),\\
\nu_\mu N &\to \nu_\mu X,\qquad (X~\text{hadrons, no muon}),
\end{align}
which are second order in the Fermi theory and hence suppressed. Nevertheless, the experiment found
\begin{equation}
R_\nu=\frac{\sigma_{\rm nc}(\nu_\mu N\to \nu_\mu X)}{\sigma_{\rm cc}(\nu_\mu N\to \mu^- X)}=0.31\pm 0.01,
\qquad
R_{\bar\nu}=\frac{\sigma_{\rm nc}(\bar\nu_\mu N\to \bar\nu_\mu X)}{\sigma_{\rm cc}(\bar\nu_\mu N\to \mu^+ X)}=0.38\pm 0.02,
\end{equation}
showing that charged and neutral currents have roughly the same size.

\section{Parametrizing neutral currents and extracting couplings from DIS}
Neutral currents are parameterized in the slides as
\begin{equation}
J^\mu_{\rm nc}=\sum_f \bar f\,\gamma^\mu\,\frac{c_V^f-c_A^f\gamma^5}{2}\,f,
\end{equation}
and the corresponding interaction is
\begin{equation}
\boxed{
\mathcal{L}_{\rm nc}
=
4\sqrt{2}\,G\,\rho\;J^{\mu\dagger}_{\rm nc}J_{{\rm nc}\,\mu}.
}
\end{equation}
Here $f$ runs over fermion fields (quarks and leptons). The parameters $\rho$, $c_V^f$, $c_A^f$
can be measured in high-energy neutrino deep inelastic scattering (DIS).

\subsection{DIS kinematics and the parton model benchmark}
Processes:
\begin{align}
\nu_\ell N &\to \nu_\ell X,\qquad &\nu_\ell N &\to \ell^- X,\\
\bar\nu_\ell N &\to \bar\nu_\ell X,\qquad &\bar\nu_\ell N &\to \ell^+ X,
\qquad \ell=e,\mu.
\end{align}
For electron DIS, the parton model gives (for $s\gg m_N$)
\begin{equation}
\frac{d\sigma}{dx\,dy}
\simeq
F_2(x)\,\frac{2\pi\alpha^2}{s}\,\frac{1}{q^4}\left(1+(1-y)^2\right),
\qquad
F_2(x)=\sum_i f_i(x)Q_i^2\,x.
\end{equation}
Kinematic variables as defined in the slides:
\begin{equation}
x=-\frac{q^2}{2m_N\nu},
\qquad
\nu=\frac{p_N\cdot q}{m_N},
\qquad
y=\frac{p_N\cdot q}{p_N\cdot p_A},
\qquad
q=p_A-p_1,
\end{equation}
where $p_A$ is the incoming electron momentum and $p_1$ the outgoing electron momentum.
From data, one extracts $F_2(x)$ and thus combinations of parton distributions (e.g.\ $xf_{uv}(x)$, $xf_{dv}(x)$, $xf_s(x)$).

\subsection{Charged-current neutrino DIS formulas}
Analogous formulas in the slides for charged-current neutrino DIS:
\begin{align}
\frac{d\sigma(\nu_\ell N\to \ell^- X)}{dx\,dy}
&=
\frac{G^2 s}{\pi}\left(x f_d(x) + x f_{\bar u}(x)(1-y)^2\right),\\
\frac{d\sigma(\bar\nu_\ell N\to \ell^+ X)}{dx\,dy}
&=
\frac{G^2 s}{\pi}\left(x f_{\bar d}(x) + x f_{u}(x)(1-y)^2\right).
\end{align}
From these one can extract $xf_{uv}(x)$, $xf_{dv}(x)$ and $xf_s(x)$, i.e.\ how nucleons are made of partons.

\subsection{Neutral-current DIS and extracted couplings}
Similar formulas for
$\nu_\ell N\to \nu_\ell X$ and $\bar\nu_\ell N\to \bar\nu_\ell X$
allow extraction of $\rho$, $c_V^f$, $c_A^f$ and the parton charges $Q_f$.
The slides report:
\begin{equation}
\rho\simeq 1,
\end{equation}
and the following values:
\begin{center}
\begin{tabular}{@{}cccc@{}}
\toprule
$f$ & $c_V^f$ & $c_A^f$ & $Q_f$\\
\midrule
$\nu_\ell$ & $1/2$ & $1/2$ & $0$\\
$e^-$ & $-0.03$ & $-1/2$ & $-1$\\
$u$ & $0.19$ & $1/2$ & $2/3$\\
$d$ & $-0.34$ & $-1/2$ & $-1/3$\\
\bottomrule
\end{tabular}
\end{center}
A remarkable relation holds:
\begin{equation}
\boxed{
c_V^f = c_A^f - 2Q_f\,x,
\qquad x\simeq 0.23.
}
\end{equation}
Therefore the neutral current can be written (as in the slides) as
\begin{equation}
\boxed{
J^\mu_{\rm nc}
=
\sum_F
\bar F\left(\frac{\tau^3}{2}\gamma^\mu P_L - x\,\bar F Q_F\gamma^\mu F\right),
}
\end{equation}
where $Q_F$ is the charge matrix in each doublet:
\begin{equation}
Q_F=
\begin{pmatrix}
0&0\\
0&-1
\end{pmatrix}
\quad\text{for leptons},
\qquad
Q_F=
\begin{pmatrix}
\frac{2}{3}&0\\
0&-\frac{1}{3}
\end{pmatrix}
\quad\text{for quarks}.
\end{equation}

\section{Big picture summary (within the slides' scope)}
\begin{itemize}
  \item Weak interactions are required by the extreme smallness of certain decay widths
        (pion, neutron, muon, strange hadrons) compared to typical strong widths.
  \item The Wu experiment established \textbf{parity violation}, motivating the \textbf{V--A} structure.
  \item The Fermi constant has dimension $-2$, indicating an effective, low-energy description
        consistent with massive vector exchange ($G\sim 1/m_W^2$).
  \item Low-energy processes:
        \begin{itemize}
          \item $\pi^+\to \ell^+\nu_\ell$ exhibits \textbf{helicity suppression} ($\mathcal{M}\propto m_\ell$),
                yielding a parameter-free $e/\mu$ ratio.
          \item Chiral symmetry and the non-linear sigma model allow a systematic hadronic representation
                of weak currents, giving $F_\pi=f_\pi$ and predicting other low-energy weak processes
                without new parameters at leading order.
          \item Nucleon beta decay requires an axial coupling $g_A\simeq 1.27$ and leads, in the non-relativistic limit,
                to Fermi and Gamow--Teller operators.
        \end{itemize}
  \item Muon decay provides an independent determination of $G$ and motivates Cabibbo mixing to reconcile
        with beta decay determinations.
  \item Strange decays require different coupling strengths ($V_{us}$ vs.\ $V_{ud}$), leading to Cabibbo mixing and,
        with three families, to the CKM phase and CP violation.
  \item Neutral currents are required by an underlying $SU(2)$ structure and were discovered in DIS;
        their couplings satisfy the relation $c_V^f=c_A^f-2Q_f x$ with $x\simeq 0.23$.
  \item Neutrino masses and mixing lead to the PMNS matrix and neutrino oscillations, with possible additional
        CP-violating phases in the Majorana case.
\end{itemize}